%%%%%%%%%%  scan idx 453  |  folio 442  %%%%%%%%%%

\medskip
\noindent
10.3. \underline{Cambio de rasgo}. Sea
\begin{equation}\tag{10.3.1}
g:S'\longrightarrow S
\end{equation}
un morfismo dominante de rasgos henselianos, asociado a un
homomorfismo local inyectivo de anillos de valoración discreta
\begin{equation*}
V\longrightarrow V'.
\end{equation*}
Como de costumbre, pongamos
\begin{equation}\tag{10.3.2}
e(S'/S)=\operatorname{length}_V\bigl(V'/\underline mV'\bigr),
\end{equation}
donde $\underline m$ es el ideal maximal de $V$. Se sabe entonces que
el homomorfismo inducido por (10.3.1) sobre las partes ``moderadas'' de
los grupos de inercia,
\begin{equation}\tag{10.3.3}
g_*:I'_t\simeq\prod_{\ell\ne p}\mathbb Z_\ell(1)(k')
 \longrightarrow
 I_t=\prod_{\ell\ne p}\mathbb Z_\ell(1)(k),
\end{equation}
está dado por
\begin{equation}\tag{10.3.4}
g_*(x)=e(S'/S)x,
\end{equation}
después de identificar los dos miembros de
% Canon: SGA-SRC-001287.
\sisi{(10.3.3)}{(10.3.3)} mediante los isomorfismos evidentes
$\mathbb Z_\ell(1)(k)\xrightarrow{\ \sim\ }
\mathbb Z_\ell(1)(k')$.

Por otra parte, si $A_K$ es un esquema abeliano sobre $K$ con
reducción semiestable sobre $S$, la formación de los retículos
retorcidos $\underline M$ y $\underline M'$ conmuta con todo cambio de
base $S'\longrightarrow S$ como el anterior, según se deduce
inmediatamente de las definiciones (en especial, de 3.4). Podemos,
pues, considerar el acoplamiento $u_\ell^{A_{K'}}$ definido por
$A_{K'}$, donde $K'$ es el cuerpo de fracciones de $V'$, como un
acoplamiento
\begin{equation*}
u_\ell^{A_{K'}}:
 \underline M_{\ell S'}\otimes\underline M'_{\ell S'}
 \longrightarrow\mathbb Z_{\ell S'},
\end{equation*}
y compararlo con el acoplamiento $u_{\ell S'}^{A_K}$ obtenido de
$u_\ell^{A_K}$ por cambio de base. Cuando $\ell\ne p$, la construcción
de los acoplamientos $u_\ell$ dada en 9.2, junto con la fórmula
(10.3.4), muestra inmediatamente que
\begin{equation}\tag{10.3.5}
u_\ell^{A_{K'}}=e(S'/S)u_\ell^{A_K}.
\end{equation}

%%%%%%%%%%  scan idx 454  |  folio 443  %%%%%%%%%%

En realidad, esta fórmula sigue siendo válida para $\ell=p$. Para
verlo, volvamos a la construcción general de 9.3--9.6 y reduzcamos a
la afirmación de compatibilidad análoga para el isomorfismo 9.4.4(e),
bajo un cambio de rasgos estrictamente locales
$S'\longrightarrow S$. Obsérvese que el isomorfismo en cuestión está
inducido por la valoración normalizada $K^*/V^*\simeq\mathbb Z$ y da,
por tanto, el diagrama conmutativo
\begin{center}
\begin{tikzcd}
K^*/V^* \arrow[r,"\sim"] \arrow[d]
  & \mathbb Z \arrow[d,"e(S'/S)"] \\
K^{\prime *}/V^{\prime *} \arrow[r,"\sim"]
  & \mathbb Z .
\end{tikzcd}
\end{center}

He aquí un resultado más profundo sobre los acoplamientos de
monodromía. Lo demostraremos mediante métodos trascendentes en la
sección 12, apoyándonos en resultados que se probarán más adelante en
el Seminario. (En realidad, no utilizaremos 10.4 en el resto del
Seminario, como tampoco utilizaremos los demás resultados del presente
exposé.)

\medskip
\noindent
\underline{Teorema} 10.4. \underline{Sea $S$ un rasgo henseliano, y
sea $A_K$ un esquema abeliano sobre $K$ con reducción semiestable sobre
$S$. Consideremos los retículos $\underline M$ y $\underline M'$
asociados como en 9.1 (cf. 5.5.2). Entonces se cumplen las afirmaciones
siguientes:}

%%%%%%%%%%  scan idx 455  |  folio 444  %%%%%%%%%%

\smallskip
\noindent
a) \underline{Existe un único acoplamiento}
\begin{equation}\tag{10.4.1}
u:\underline M\otimes\underline M'\longrightarrow\mathbb Z_S
\end{equation}
\underline{tal que, para todo número primo $\ell$,
$u\otimes_{\mathbb Z}\mathbb Z_\ell$ sea el acoplamiento de monodromía
$u_\ell$ de (9.1.2). Este acoplamiento es separante.}

\smallskip
\noindent
b) \underline{Sea $\xi$ una polarización de $A_K$, que define por
tanto un homomorfismo (10.2.3)}
\begin{equation*}
\xi_t:\underline M\longrightarrow\underline M',
\end{equation*}
\underline{y consideremos la forma}
\begin{equation}\tag{10.4.2}
\begin{aligned}
\varphi_\xi^t:\underline M\otimes\underline M
  &\longrightarrow\mathbb Z_S,\\
\varphi_\xi^t(x,y)&=u\bigl(x,\xi_t(y)\bigr).
\end{aligned}
\end{equation}
\underline{Esta forma es simétrica, positiva y no degenerada.}

Naturalmente, en términos de la fibra geométrica $M_o$ de
$\underline M$ en un punto geométrico dado de $S$, la afirmación
precedente dice que la forma correspondiente
\begin{equation*}
M_o\otimes M_o\longrightarrow\mathbb Z
\end{equation*}
es simétrica, positiva y no degenerada. El acoplamiento (10.4.1) se
llamará de nuevo \underline{acoplamiento de monodromía}.
